\section{Foundations}\label{foundations}

The formal objects that carry the mathematical-physics content are a positive degree-of-freedom count, a canonical binary decision encoding, structural rank, and decision entropy.

\subsection{Formal Object}

\begin{definition}[Bounded Decision System]\label{def:architecture}
A \emph{bounded decision system} is a finite bounded physical system equipped with a positive integer $\mathrm{DOF}(A)$. The corresponding Lean object is named \texttt{Architecture}. The results used below depend only on the degree-of-freedom count and on the canonical decision encoding.
\end{definition}

\textbf{Interpretation.} $\mathrm{DOF}(A)$ counts independent coordinates that can vary separately. Subsequent sections study what that count forces once one asks for exact resolution.

\subsection{Degrees of Freedom}

\begin{definition}[Degrees of Freedom]\label{def:dof}
The quantity $\mathrm{DOF}(A) \in \mathbb{N}$ counts independent coordinates of variation in a bounded decision system $A$. In the mechanized development it is the local structural parameter attached to \texttt{Architecture}; later sections identify it exactly with the structural rank of a canonical decision problem.
\end{definition}

\textbf{Operational meaning.} If $\mathrm{DOF}(A)=n$, the system has $n$ independent coordinates that must be resolved in the worst case by any exact resolver.

\begin{proposition}[DOF Additivity]\label{prop:dof-additive}
For disjoint bounded decision systems $A_1$ and $A_2$:
\[
\mathrm{DOF}(A_1 \oplus A_2) = \mathrm{DOF}(A_1) + \mathrm{DOF}(A_2)
\]
\end{proposition}
\leanmeta{\LH{L17}, \LH{L19}}

\begin{proof}
Independent coordinate sets combine by disjoint union, so the coordinate count is additive under composition.
\end{proof}

\subsection{Finite Physical Acquisition}

\begin{theorem}[Counting Gap]\label{thm:counting-gap}
Let $\varepsilon, C \in \mathbb{N}$ with $\varepsilon>0$ and $C>0$. If each information-acquisition event consumes $\varepsilon$ discrete cost units, then
\[
\varepsilon \cdot N \le C \implies N \le C.
\]
Equivalently, any bounded system with positive per-event cost admits only finitely many acquisition events.
\end{theorem}
\leanmeta{\LH{BA10}}

\begin{proof}
In $\mathbb{N}$ every positive integer is at least one, so $N = 1\cdot N \le \varepsilon N \le C$.
\end{proof}

Theorem~\ref{thm:counting-gap} fixes the finite-event statement of the model. Proposition~\ref{prop:bounded-region} and Theorem~\ref{thm:bounded-acquisition} fix the geometric acquisition bound. Theorems~\ref{thm:discrete-acquisition} and \ref{thm:one-transition-one-bit} fix the acquisition-event interface. Landauer calibration is applied only after that interface is fixed.

\begin{proposition}[Bounded Region]\label{prop:bounded-region}
A bounded physical region is characterized by diameter $d>0$ and signal speed $c>0$. Its maximum information-acquisition rate is $c/d$ events per unit time.
\end{proposition}
\leanmeta{\LH{BA1}}

\begin{theorem}[Bounded Acquisition Rate]\label{thm:bounded-acquisition}
For a bounded region with diameter $d$, signal speed $c$, and operating time $T$,
\[
\mathrm{acquisitions}(T) \le \frac{cT}{d}.
\]
In particular, acquisition count is finite on finite horizons.
\end{theorem}
\leanmeta{\LH{BA2}}

\begin{proof}
Signals require at least $d/c$ time to traverse the region, so no more than $c/d$ acquisition events can occur per unit time.
\end{proof}

\begin{theorem}[Discrete Acquisition]\label{thm:discrete-acquisition}
In the imported bounded-acquisition model, information acquisition is counted by transition points of a finite discrete system. Acquisition counts are therefore discrete event counts.
\end{theorem}
\leanmeta{\LH{BA3}}

\begin{proof}
The imported model represents a bounded physical decision process by a finite \texttt{DiscreteSystem}. Its acquisition count is \texttt{bitOperations}, which counts transition points along a run.
\end{proof}

\begin{theorem}[One Transition, One Bit]\label{thm:one-transition-one-bit}
In the imported discrete acquisition model, each transition point contributes one unit to the bit-operation count. The canonical binary encoding therefore uses one Boolean coordinate per elementary acquisition event.
\end{theorem}
\leanmeta{\LH{BA4}}

\begin{proof}
The imported theorem states that a transition point at time $t$ contributes at least one unit to the acquisition count up to time $t+1$. The model therefore treats each elementary transition as one Boolean acquisition event.
\end{proof}

Together, Theorems~\ref{thm:bounded-acquisition}, \ref{thm:discrete-acquisition}, and \ref{thm:one-transition-one-bit} identify the later binary decision encoding as the natural finite acquisition model. A bounded resolver acquires information through finitely many discrete events, and each such event contributes one elementary boolean distinction.

\begin{theorem}[Resolution Requires a Sufficient Coordinate Set]\label{thm:resolution-sufficient}
Any exact physical resolver for a decision problem must read a sufficient coordinate set. If fewer coordinates are read, there exist states indistinguishable to the resolver but requiring different optimal actions.
\end{theorem}
\leanmeta{\LH{BA5}}

\begin{proof}
If the accessed coordinates are not sufficient, two states agree on every read coordinate while disagreeing on the optimal action. Any resolver limited to those coordinates must therefore err on at least one of the two states.
\end{proof}

\subsection{Canonical Decision Encoding}

\begin{definition}[Canonical Decision Problem]\label{def:canonical-dp}
For a bounded decision system $A$ with $\mathrm{DOF}(A)=n$, the canonical decision problem
\[
\mathrm{canonicalDP}(A)
\]
has state space $\mathrm{Fin}\;n \to \mathrm{Bool}$ and action space $\mathrm{Fin}\;n \oplus \mathrm{Unit}$. Action $\mathrm{inl}(i)$ queries coordinate $i$; the fallback action $\mathrm{inr}(\star)$ receives constant utility $1$. Query action $i$ receives utility $2$ exactly when coordinate $i$ is true and $0$ otherwise.
\end{definition}

The encoding is the exact Lean object \texttt{canonicalDP} in \texttt{Leverage/BridgeToDQ.lean}. The paper studies this encoding as the exact finite-resolution object attached to the declared degree-of-freedom count. It assigns one binary acquisition channel to each degree of freedom and one query action to each channel. The next section identifies the structural rank of this object with that coordinate count.

\subsection{Structural Rank}

\begin{definition}[Structural Rank]\label{def:srank}
The \emph{structural rank} of a finite decision problem is the cardinality of its relevant coordinate set, equivalently the size of any minimal sufficient set. It is the minimum interaction dimension that must be read to determine the optimal action exactly.
\end{definition}

For the canonical decision problem attached to $A$, the relevant coordinate set is all of $\mathrm{Fin}\;\mathrm{DOF}(A)$, so the structural-rank problem is exactly matched to the local degree-of-freedom count.

\subsection{Decision Quotient and Entropy}

\begin{definition}[Decision Quotient]\label{def:decision-quotient}
For a decision problem $D$, states are identified when they induce the same optimal-action set. The quotient space of these equivalence classes is the \emph{decision quotient} of $D$.
\end{definition}

\begin{proposition}[Optimizer Quotient as Canonical Exact Abstraction]\label{prop:optimizer-quotient}
Let $\Opt : S \to \mathcal P(A)$ be the optimizer map of a finite decision problem. In \textbf{Set}, the decision quotient is the coimage of $\Opt$, canonically equivalent to $\operatorname{im}(\Opt)$ \cite{maclane1998categories}. Any surjective decision-preserving summary factors through this quotient.
\end{proposition}
\leanmeta{\LHrng{QT}{1}{3}, \LH{QT5}, \LH{QT7}}

The quotient is the coarsest exact abstraction of the decision problem: it forgets only decision-irrelevant distinctions and preserves every distinction needed for exact action selection. The entropy and thermodynamic bounds below are stated for this canonical exact abstraction.

\begin{definition}[Decision Entropy]\label{def:decision-entropy}
Let $\mathrm{numOptClasses}(D)$ be the number of equivalence classes in the decision quotient. Two entropy normalizations are used:
\[
H_{\mathrm{bits}}(D) = \log_2 \mathrm{numOptClasses}(D),
\qquad
H_{\mathrm{nats}}(D) = \log \mathrm{numOptClasses}(D).
\]
\end{definition}

The physics results are naturally stated in nats because Landauer calibration contributes the factor $k_B T$ per nat of resolved decision information \cite{landauer1961irreversibility,bennett1982thermodynamics}.

\subsection{Formalization in Lean}

The local degree-of-freedom object lives in \texttt{Leverage/Foundations.lean}, while the canonical decision encoding and its rank-identification theorems live in \texttt{Leverage/BridgeToDQ.lean}. Structural rank and decision entropy are formalized in the decision-quotient development.
